\subsection{Квоты и снапшоты}

\subsubsection{Квота}

Квота позволяет ограничить объём дискового пространства, доступного для использования во всём томе или в отдельном каталоге. Квота настроена на томе \code{volume1}: включена и установлен лимит 10~МБ на корень тома (листинг~\ref{lst:quota1}).
\begin{Verbatim}
sudo gluster volume quota volume1 enable
sudo gluster volume quota volume1 limit-usage / 10MB
sudo gluster volume quota volume1 list
\end{Verbatim}

Первая проверка показала особенность реализации: утилита \code{dd} записала 20~МБ без ошибки, а превышение было обнаружено только после пересчёта занятого места. Квота проверяется не при каждой записи, а с задержкой, определяемой параметрами \code{features.soft-timeout} и \code{features.hard-timeout}. После этого создание нового каталога уже блокировалось ошибкой «Превышена дисковая квота».

\lstfile{6_1_quota_default_timeout.txt}{Квота на том при настройках задержки по умолчанию}\label{lst:quota1}

Для немедленного срабатывания квоты оба параметра заданы равными нулю. Повторная запись 20~МБ остановилась на отметке около 10~МБ с ошибкой \code{Disk quota exceeded} (небольшое превышение лимита связано с блочным учётом). Дополнительно установлена квота 2~МБ на каталог \code{/dir1}: запись в него прекратилась после превышения этого лимита, хотя до лимита всего тома ещё оставалось свободное место (листинг~\ref{lst:quota2}). После проверки квота на каталог удалена, квота на корень тома оставлена включённой.

\lstfile{6_1_quota_zero_timeout.txt}{Срабатывание квот на том и на каталог}\label{lst:quota2}

\subsubsection{Снапшот}

Снапшот~--- мгновенный снимок состояния тома GlusterFS на определённый момент времени. Он создаётся средствами LVM и не копирует данные физически, поэтому не требует остановки тома. GlusterFS поддерживает снапшоты только для бриков, расположенных на logical volume с thin-provisioning, поэтому брики обычных каталогов для этой цели не подходят.

Для снапшотов выбран том \code{volume2}. На \code{node1} и \code{node2} в VirtualBox добавлено по одному виртуальному диску объёмом 10~ГБ (\code{/dev/sdb}). На каждом из узлов создан физический том, группа томов \code{vg_brick2}, thin-пул \code{pool2} объёмом 8~ГБ и thin-том \code{lv_brick2} объёмом 6~ГБ, отформатированный в XFS и смонтированный в \code{/bricks/brick2}. Брик расположен в подкаталоге \code{/bricks/brick2/brick}, чтобы данные не попали в корневую файловую систему, если том по какой-либо причине не смонтируется. Старый \code{volume2} удалён и создан заново на LVM-бриках (листинг~\ref{lst:lvm}).

\lstfile{6_2_lvm_bricks.txt}{Подготовка LVM-бриков и пересоздание тома volume2}\label{lst:lvm}

Снапшот создан командой \code{gluster snapshot create}. Ключ \code{no-timestamp} убирает из имени метку времени. Снапшот создаётся в неактивном состоянии (в статусе \code{Deactivated Snapshot}), на каждом брике появляется LVM-снимок, у которого в качестве исходного тома указан \code{lv_brick2}. Затем в томе удалены и изменены файлы, а том остановлен и восстановлен из снапшота командой \code{snapshot restore}. После восстановления и повторного монтирования данные вернулись к состоянию на момент создания снапшота. Использованный при восстановлении снапшот расходуется и исчезает из списка. Дополнительно создан и удалён второй снапшот (листинг~\ref{lst:snap}).

\lstfile{6_3_snapshot.txt}{Создание снапшота, восстановление данных и удаление снапшота}\label{lst:snap}
